\input{../tex-inputs/latex-leseansicht-vorspann}

\section[Arthur Schnitzler an Gustav Schwarzkopf, 22. 8. 1923]{L04350 Arthur Schnitzler an Gustav Schwarzkopf, 22. 8. 1923}
\nopagebreak\mylabel{L04350v}
\rehead{ }\normalsize\beginnumbering\briefempfaengerindex{Schwarzkopf, Gustav@\textsc{Schwarzkopf, Gustav}!zzzSchnitzler, Arthur@\emph{von Arthur Schnitzler}!1923-08-221@{22. 8. 1923}|(be}
\toendnotes[C]{\smallbreak\pagebreak[2]}
\correspDesc{Versand  durch Arthur Schnitzler am 22. 8. 1923 in Celerina
\newline{}Übermittlung  am 22. 8. 1923 in Celerina
\newline{}Erhalt  durch Gustav Schwarzkopf im Zeitraum [23. 8. 1923 – 27. 8. 1923?] in Wien}\toendnotes[C]{\smallbreak}
\Standort{DLA, A:Schnitzler, HS.1985.1.1897.}
\physDesc{Bildpostkarte, 558 Zeichen
\newline{}Handschrift: Bleistift, lateinische Kurrent
\newline{}Versand: Stempel: »\nobreak{}\oindex{Celerina@\textbf{Celerina}|pwk}Celerina (Graubünden), 22 VIII 23, 15\nobreak{}«.  }\pstart{}{\pb}Hr Gustav Schwarzkopf\pend{}\pstart{}Wien I\oindex{I., Innere Stadt@\textbf{I., Innere Stadt}, \emph{Verwaltungsgebiet}|pw}.\pend{}\pstart{}Tiefer Graben 17\oindex{Wien@\textbf{Wien}!I., Innere Stadt@\textbf{I., Innere Stadt}!Tiefer Graben 17@\textbf{Tiefer Graben 17}, \emph{Wohngebäude}|pw}.\pend{}{\bigskip}
\pstart
           \noindent{}\centering{}{\pb}\textcolor{gray}{\textbf{Celerina\oindex{Celerina@\textbf{Celerina}|pw} mit Berninabahn\oindex{Berninabahn@\textbf{Berninabahn}, \emph{Bahn}|pw}}}\pend
           \vspace{1em}
\pstart
           \raggedleft{}{\pb}Celerina, Cresta Palace\oindex{Cresta Palace@\textbf{Cresta Palace}, \emph{Hotel}|pw}\pend
           
\pstart
           \raggedleft{}22. 8. 923\pend
           \vspace{0.5em}
\pstart
           lieber Guſtav, seien Sie sehr herzlichst gegrüßt!  Ich war in Salzburg\oindex{Salzburg@\textbf{Salzburg}, \emph{Verwaltungsgebiet}|pw}, Baden-Baden\oindex{Baden-Baden@\textbf{Baden-Baden}|pw}, am Bodensee\oindex{Bodensee@\textbf{Bodensee}, \emph{See}|pw} u in Bayern\oindex{Bayern@\textbf{Bayern}, \emph{Land}|pw}, eh ich hier in diesem wunderbaren Ort \introOben{}gelandet bin u\introOben{} mich von der köstlichen Luft und den kleinen
               Ziffern berauschen lasse. Wie wohl thut unter solchen Umständen eine à conto Zahlung
               von Fischer\pwindex{Fischer, Samuel 24.\,12.\,1859 Liptovský Mikuláš – 15.\,10.\,1934 Berlin@\textsc{Fischer, Samuel} (24.\,12.\,1859 Liptovský Mikuláš – 15.\,10.\,1934 Berlin), \emph{Verleger}|pw} in Berlin\oindex{Berlin@\textbf{Berlin}, \emph{Hauptstadt}|pw} von 10 Millionen Mark = 40 francs! (für {\pb}die man in Wirklichkeit nach Qualität und  Quantität dort mehr kriegt als für 10
               Mill Mark)\pend
           
\pstart
           Auch an Dr.  Max\pwindex{Schwarzkopf, Max 12.\,6.\,1857 Wien – 14.\,4.\,1928 ebd.@\textsc{Schwarzkopf, Max} (12.\,6.\,1857 Wien – 14.\,4.\,1928 ebd.), \emph{Rechtsanwalt}|pw} alles herzliche{\\[\baselineskip]}Ihr getreuer{\\[\baselineskip]}\spacefill\mbox{A. S.}\pend
           \leftskip=0em{}\selectlanguage{ngerman}\endnumbering\briefempfaengerindex{Schwarzkopf, Gustav@\textsc{Schwarzkopf, Gustav}!zzzSchnitzler, Arthur@\emph{von Arthur Schnitzler}!1923-08-221@{22. 8. 1923}|)be}\mylabel{L04350h}
\begin{anhang}
\end{anhang}\newcommand{\dateiname}{L04350}\newcommand{\titel}{Arthur Schnitzler an Gustav Schwarzkopf, 22. 8. 1923}\newcommand{\editorInnen}{Herausgegeben von Jahnke, SelmaMüller, Martin Anton}\input{../tex-inputs/latex-leseansicht-abspann}
